\Chapitre{Normes de dessins des schémas}

\section{Généralité}

Ces normes doivent être appliquées à tous les nouveaux schémas et dans la mesure du possible aux modifications de schémas déjà dessinés.

\section{Principes}

Tous câblage doit avoir son schéma.

Les ensembles standardisés peuvent avoir un seul schéma de principe commun pour une même version d'ensemble.

Tous borniers, connecteurs et entrée/sortie d'automate doit avoir son affectation.

Tout relais doit avoir son affectation.

Toute voie doit apparaitre sur un plan de voie.

Les schémas et plan de voie doivent être dessiné avant le câblage.

\section{Numérotation}

\subsection{Cartouche}

Le cartouche doit contenir au moins : le numéro du document, le nom du dessinateur initial, la date du dessin initial, le nom du dessinateur de la dernière modification, la date de la dernière modification et le nom de la dernière modification.

\subsection{Format}

La numérotation des schémas suit le format suivant :

AAA-000-012

AAA : abréviation de l'enclenchement/du satellite ou nom du module de pleine voie.\\
000 : contenu du schéma selon la liste ci-dessous.\\
123 : numérotation continue (le premier chiffre peut également servir à désigner le contenu)

Exemple : SOU-604-003\\
Un schéma de commande d'itinéraire de l'enclenchement de Souche, page 3.
